\documentclass[10pt,a4paper]{article}
\usepackage[T1]{fontenc}
\usepackage[utf8]{inputenc}
\usepackage[ngerman]{babel}
\usepackage[a4paper,left=12mm,right=12mm,top=12mm,bottom=9mm]{geometry}
\usepackage{lmodern}
\usepackage[table]{xcolor}
\usepackage{tabularx}
\usepackage{array}
\usepackage{microtype}
\usepackage[most]{tcolorbox}
\pagestyle{empty}
\setlength{\parindent}{0pt}
\setlength{\parskip}{0pt}
\setlength{\tabcolsep}{3pt}
\renewcommand{\familydefault}{\sfdefault}
\linespread{1.0}\selectfont
\renewcommand{\arraystretch}{1.02}
\definecolor{Navy}{HTML}{17324D}
\definecolor{Blue}{HTML}{3B6EA5}
\definecolor{Sky}{HTML}{EAF3FA}
\definecolor{Mint}{HTML}{E8F5EF}
\definecolor{Gold}{HTML}{F3B33D}
\definecolor{Ink}{HTML}{1E2935}
\definecolor{Muted}{HTML}{617080}
\definecolor{Line}{HTML}{D8E1E8}
\definecolor{CodeBg}{HTML}{F5F7F9}
\color{Ink}
\newcommand{\sectiontitle}[1]{%
  \vspace{.3mm}{\color{Blue}\bfseries\large #1}\par
  \vspace{.2mm}{\color{Blue}\rule{\linewidth}{.65pt}}\par\vspace{.3mm}}
\newcommand{\sqlkw}[1]{\textcolor{Blue}{\bfseries #1}}
\newcommand{\sqlbox}[1]{%
  \begin{tcolorbox}[colback=CodeBg,colframe=Line,arc=1.3mm,boxrule=.55pt,
    left=2.2mm,right=2.2mm,top=.5mm,bottom=.5mm,boxsep=0pt]
  {\ttfamily\fontsize{9.4}{11.2}\selectfont #1}
  \end{tcolorbox}}
\newcommand{\workbox}[1]{%
  \begin{tcolorbox}[colback=Sky,colframe=Blue,arc=1.2mm,boxrule=.55pt,
    left=2mm,right=2mm,top=.35mm,bottom=.35mm,boxsep=0pt]
  {\fontsize{9.2}{11}\selectfont\hyphenpenalty=10000\exhyphenpenalty=10000\textcolor{Blue}{\bfseries Arbeitsauftrag:} #1}
  \end{tcolorbox}}
\newcommand{\pageheader}{%
  \colorbox{Navy}{\parbox{\dimexpr\linewidth-2\fboxsep\relax}{%
    \vspace{1.5mm}\color{white}\begin{tabularx}{\linewidth}{@{}Xr@{}}
    {\small\bfseries INFORMATIK \textperiodcentered{} KLASSE 9} &
    \colorbox{Gold}{\color{Navy}\small\bfseries\strut LÖSUNGEN}
    \end{tabularx}
    \vspace{1mm}{\fontsize{17}{20}\selectfont\bfseries
    Aufgabe 2 -- Erste SQL-Abfragen}\par
    \vspace{.6mm}{\large Sicherungsblatt}\vspace{1.5mm}}}}
\newcommand{\pagefooter}[1]{%
  \vfill{\color{Line}\rule{\linewidth}{.5pt}}\par\vspace{1mm}
  \begin{tabularx}{\linewidth}{@{}Xc r@{}}
  {\small\color{Muted}Klasse 9 \textperiodcentered{} Datenbanken} &
  {\small\color{Muted}#1 / 2} &
  {\small\color{Muted}Sicherungsblatt \textperiodcentered{} Aufgabe 2}
  \end{tabularx}}
\newcommand{\merkebox}[1]{%
  \begin{tcolorbox}[colback=Mint,colframe=Blue,arc=2mm,boxrule=.8pt,
    left=3mm,right=3mm,top=1.5mm,bottom=1.5mm]
  \textcolor{Blue}{\bfseries Merke}\par\vspace{.5mm}
  {\fontsize{9.6}{11.6}\selectfont #1}
  \end{tcolorbox}}
\newcommand{\smallcaption}[1]{{\small\color{Muted}#1}\par\vspace{.5mm}}
\newcommand{\miniresult}[2]{%
  {\fontsize{8.6}{10.2}\selectfont
  \begin{tabular}{|>{\ttfamily\bfseries}l|}\hline
  \rowcolor{Sky}#1\\\hline
  #2\\\hline
  \end{tabular}}}
\begin{document}
\pageheader\par
\sectiontitle{1 \textperiodcentered{} Von users zur Ergebnisrelation}
\smallcaption{Ausschnitt aus users (4 von 317 Zeilen; 5 von 18 Spalten)}
{\fontsize{8.6}{10.2}\selectfont
\begin{tabularx}{\linewidth}{|>{\ttfamily\bfseries}p{7mm}|>{\raggedright\arraybackslash}X|>{\ttfamily\raggedright\arraybackslash}p{31mm}|>{\ttfamily\raggedright\arraybackslash}p{25mm}|>{\raggedright\arraybackslash}p{19mm}|}\hline
\rowcolor{Sky}id & name & username & birthday & city\\\hline
1 & Noah Schneider & zornigerfotograf & 2007-10-02 & Schwerin\\\hline
2 & Finn Becker & wanderwut\_finn & 2006-06-14 & München\\\hline
3 & Liam Hoffmann & bergcoder & 2009-07-08 & Zürich\\\hline
4 & Elias Wagner & koch\_kicker & 2008-04-23 & Berlin\\\hline
\end{tabularx}}\par
\vspace{.3mm}
\workbox{Zeige die Geburtsdaten aller Mitglieder.}
\begin{minipage}[t]{.32\linewidth}
\smallcaption{SQL-Anweisung}
\sqlbox{\sqlkw{SELECT} birthday\par\sqlkw{FROM} users}
\end{minipage}\hfill
\begin{minipage}[t]{.64\linewidth}
\smallcaption{Ergebnisrelation}
\miniresult{birthday}{2007-10-02\\2006-06-14\\2009-07-08\\2008-04-23\\\ldots}
\end{minipage}\par
\vspace{.7mm}
\smallcaption{Allgemeine Form}
\sqlbox{\sqlkw{SELECT} Spaltenname\par\sqlkw{FROM} tabellenname}

\sectiontitle{2 \textperiodcentered{} Mehrere oder alle Spalten}
\workbox{Zeige die Usernames und die Geburtsdaten aller Mitglieder an}
\begin{minipage}[t]{.27\linewidth}
\smallcaption{SQL-Anweisung}
\sqlbox{\sqlkw{SELECT} username, birthday\par\sqlkw{FROM} users}
\end{minipage}\hfill
\begin{minipage}[t]{.71\linewidth}
\smallcaption{Ergebnisrelation (erste 2 von 317 Zeilen)}
{\fontsize{8.6}{10.2}\selectfont
\begin{tabular}{|>{\ttfamily\bfseries}l|>{\ttfamily}l|}\hline
\rowcolor{Sky}username & birthday\\\hline
zornigerfotograf & 2007-10-02\\\hline
wanderwut\_finn & 2006-06-14\\\hline
\ldots & \ldots\\\hline
\end{tabular}}
\end{minipage}\par
\vspace{.2mm}
\workbox{Zeige alle Daten aller Mitglieder an.}
\begin{minipage}[t]{.22\linewidth}
\smallcaption{SQL-Anweisung}
\sqlbox{\sqlkw{SELECT} *\par\sqlkw{FROM} users}
\end{minipage}\hfill
\begin{minipage}[t]{.76\linewidth}
\smallcaption{Ergebnisrelation (2 von 317 Zeilen, 5 von 18 Spalten gezeigt)}
{\fontsize{8.6}{10.2}\selectfont\setlength{\tabcolsep}{1.5pt}
\begin{tabular}{|>{\ttfamily}c|>{\raggedright\arraybackslash}p{26mm}|>{\ttfamily\raggedright\arraybackslash}p{30mm}|>{\ttfamily}c|>{\ttfamily\raggedright\arraybackslash}p{22mm}|>{\raggedright\arraybackslash}p{18mm}|>{\ttfamily}c|}\hline
\rowcolor{Sky}id & \mbox{name} & \mbox{username} & \ldots & \mbox{birthday} & \mbox{city} & \ldots\\\hline
1 & Noah Schneider & zornigerfotograf & \ldots & 2007-10-02 & Schwerin & \ldots\\\hline
2 & Finn Becker & wanderwut\_finn & \ldots & 2006-06-14 & München & \ldots\\\hline
\ldots & \ldots & \ldots & \ldots & \ldots & \ldots & \ldots\\\hline
\end{tabular}}
\end{minipage}\par
\vspace{.2mm}
\merkebox{In der SELECT-Klausel (Fachbegriff: Projektion) werden die Spalten angegeben, die ausgegeben werden sollen. Die Spalten müssen in der Tabelle der FROM-Klausel vorhanden sein. SELECT * gibt alle Spalten aus.}

\sectiontitle{3 \textperiodcentered{} Sortieren mit ORDER BY}
\workbox{Zeige Usernames und Gr\"o\ss{}en aller Mitglieder an, nach Gr\"o\ss{}e absteigend geordnet.}
\begin{minipage}[t]{.39\linewidth}
\smallcaption{SQL-Anweisung}
\sqlbox{\sqlkw{SELECT} username, centimeters\par\sqlkw{FROM} users\par\sqlkw{ORDER BY} centimeters DESC}
\end{minipage}\hfill
\begin{minipage}[t]{.59\linewidth}
\smallcaption{Beginn der Ergebnisrelation (5 von 317 Zeilen; gleich gro\ss{}e Mitglieder können ihre Reihenfolge tauschen)}
{\fontsize{8.6}{10.2}\selectfont
\begin{tabular}{|>{\ttfamily\bfseries}l|>{\ttfamily}r|}\hline
\rowcolor{Sky}username & centimeters\\\hline
frageriese & 195\\\hline
wanderriese\_taro & 195\\\hline
turmriese\_max & 195\\\hline
reiseriese\_elias & 195\\\hline
fragezornriese & 194\\\hline
\end{tabular}}
\end{minipage}\par
\vspace{.2mm}
\merkebox{ORDER BY centimeters DESC sortiert absteigend: gro\ss{}e Werte stehen oben. Ohne DESC gilt ASC (aufsteigend).}
\pagefooter{1}

\newpage
\pageheader\par
\sectiontitle{4 \textperiodcentered{} Dopplungen vermeiden mit DISTINCT}
\begin{minipage}[t]{.48\linewidth}
\workbox{Zeige die Geschlechtsangaben aller Mitglieder an.}
\smallcaption{SQL-Anweisung}
\sqlbox{\sqlkw{SELECT} gender\par\sqlkw{FROM} users}
\smallcaption{Ergebnisrelation (Ausschnitt aus 317 Zeilen)}
{\fontsize{8.6}{10.2}\selectfont
\begin{tabular}{|>{\ttfamily\bfseries}l|}\hline
\rowcolor{Sky}gender\\\hline
male\\\hline male\\\hline male\\\hline \ldots\\\hline
\end{tabular}}
\end{minipage}\hfill
\begin{minipage}[t]{.48\linewidth}
\workbox{Zeige jede unterschiedliche Geschlechtsangabe nur einmal an.}
\smallcaption{SQL-Anweisung}
\sqlbox{\sqlkw{SELECT} \sqlkw{DISTINCT} gender\par\sqlkw{FROM} users}
\smallcaption{Ergebnisrelation (4 Zeilen)}
{\fontsize{8.6}{10.2}\selectfont
\begin{tabular}{|>{\ttfamily\bfseries}l|}\hline
\rowcolor{Sky}gender\\\hline
male\\\hline female\\\hline diverse\\\hline NULL\\\hline
\end{tabular}}\par\vspace{.5mm}
{\small\color{Muted}NULL steht für das Konto admin ohne Eintrag.}
\par\vspace{.4mm}
\smallcaption{Allgemeine Form}
\sqlbox{\sqlkw{SELECT} \sqlkw{DISTINCT} Spaltenname\par\sqlkw{FROM} tabellenname}
\end{minipage}\par
\vspace{1mm}
\merkebox{DISTINCT steht direkt hinter SELECT. Gleiche Werte erscheinen in der Ergebnisrelation nur einmal; die Tabelle users bleibt unverändert. NULL bedeutet: kein Wert eingetragen.}

\sectiontitle{5 \textperiodcentered{} Zeilen begrenzen mit LIMIT}
\begin{minipage}[t]{.48\linewidth}
\workbox{Zeige die Usernames der ersten drei Mitglieder an.}
\smallcaption{SQL-Anweisung}
\sqlbox{\sqlkw{SELECT} username\par\sqlkw{FROM} users\par\sqlkw{LIMIT} 3}
\smallcaption{Ergebnisrelation (3 Zeilen)}
\miniresult{username}{zornigerfotograf\\wanderwut\_finn\\bergcoder}
\vspace{.5mm}\par
{\fontsize{9.2}{11}\selectfont LIMIT 3 zeigt höchstens drei Zeilen; ohne Sortierung hängt die Auswahl von der vorhandenen Reihenfolge ab.}
\end{minipage}\hfill
\begin{minipage}[t]{.48\linewidth}
\workbox{Zeige die Usernames und Geburtsdaten der vier jüngsten Mitglieder an.}
\smallcaption{SQL-Anweisung}
\sqlbox{\sqlkw{SELECT} username, birthday\par\sqlkw{FROM} users\par\sqlkw{ORDER BY} birthday DESC\par\sqlkw{LIMIT} 4}
\smallcaption{Ergebnisrelation (4 Zeilen)}
{\fontsize{8.6}{10.2}\selectfont
\begin{tabular}{|>{\ttfamily\bfseries}l|>{\ttfamily}l|}\hline
\rowcolor{Sky}username & birthday\\\hline
gartenlesemaus & 2010-07-07\\\hline
gartencoderin & 2010-06-28\\\hline
codekickemmi & 2010-06-22\\\hline
modereiselio & 2010-06-15\\\hline
\end{tabular}}\par\vspace{.5mm}
{\fontsize{9.2}{11}\selectfont Erst sortiert ORDER BY die spätesten Geburtsdaten nach oben; danach zeigt LIMIT die ersten vier Zeilen.}
\end{minipage}\par
\vspace{.4mm}
\smallcaption{Allgemeine Form}
\sqlbox{\sqlkw{SELECT} Spaltenname\par\sqlkw{FROM} tabellenname\par\sqlkw{LIMIT} Anzahl}
\vspace{1mm}
\merkebox{Reihenfolge der Klauseln: SELECT (gegebenenfalls mit DISTINCT) $\rightarrow$ FROM $\rightarrow$ ORDER BY $\rightarrow$ LIMIT.}
\pagefooter{2}
\end{document}
